\documentclass[10pt]{article}

\usepackage[utf8]{inputenc}

\usepackage[T1]{fontenc}

\usepackage{amsmath}

\usepackage{amsfonts}

\usepackage{amssymb}

\usepackage[version=4]{mhchem}

\usepackage{stmaryrd}

\usepackage{mathrsfs}



\begin{document}

пространство пробных функций, $\mathfrak{D}-$ гильбертово пространство, являющееся пополнением $V$ по выделенному в нем скалярному произведению, а $\mathscr{V}$ ' - топологически двойственное пространство непрерывных линейных функционалов на $\mathscr{V}$ Пусть $\mu$ есть $\mathcal{V}$-квазиинвариантная мера на $\mathscr{V}$ '. Тогда в гильбертовом пространстве $H=L_{m}^{2}\left(\mathscr{V}^{\prime}\right)$ действует представление вейлевской формы перестановочных соотношений вида



\[

\left.\begin{array}{c}

(U(f) \psi)(F)=e^{i F(f)} \psi(F), \\

(V(g) \psi)(F)=c(g, F) \sqrt{\frac{d \mu(F+g)}{d \mu(F)}} \psi(F+g), \\

\psi(F) \in H,

\end{array}\right\}

\]



где $c(g, F)$ удовлетворяет условию коцикла



\[

c(g, F) c(h, F+g)=c(g+h, F) .

\]



Это представление является циклическим относительно операторов $U(f)$ и имеет циклич. вектором функцию $\Omega(F)=1, \mu-$ П. В. $F$. В заданном представлении теорема Стоуна позволяет восстановить генераторы, реализующие форму Гейзенберга перестановочных соотношений. Представлению Фока отвечает гауссовская мера. Это единственное представление, обладающее универсально инвариантным вакуумом. Эргодичность $\mu$ является достаточным условием неприводимости представления. Два представления ун и тарн о эк в и вале нтны, если соответствующие им меры абсолютно непрерывны, а коциклы $c(g, F)$ и $c^{\prime}(g, F)$ когомологичны, то есть существует функция $K(F)$ такая, что



\[

c^{\prime}(g, F)=K(F) c(g, F) K(F)^{-1} .

\]



Порождающий функщионал



\[

E(f)=(\Omega(F), U(f) \Omega(F))=\int e^{i F(f)} d \mu(F)

\]



меры $\mu$ определяет (регулярное) состояние на $C^{*}$-алгебре, порождаемой операторами



\[

\{U(f) ; V(f)\}

\]



что дает связь схемы Гельфанда - Араки со схемой, основанной на ГНС-конструкции.



Jum.: [1] Garding L., Wightman A., «Proc. Nat. Acad. Sci. USA», 1954, v. 70, № 7, p. 617-25; [2] Wight man A. S., Schwe ber S.S., "Phys. Rev.», 1955, v. 98, № 3, p. 812-37; [3] Голодец В.Я., «Успехи математич. наук», 1969, т. 24, в. 4, с. 3-64; [4] KlauderJ.R., MckennaJ., Woods E.J., UJ. Math. Phys.», 1966, v. 7, № 5, p. 822-28; [5] ГельфандИ. М., Виленк и н Н. Я., Некоторые применения гармонического анализа: оснащенные гильбертовы пространства, М., 1961; [6] Araki H., «J. Math. Phys.», 1960, v. 1, № 6, p. 492-504; [7] Fu ku to me H., «Progr. Theor. Phys.», 1960, v. 23, № 6, p. 989-1002; [8] С и гал И., Математические проблемы релятивистской физики, пер. с англ., М., $\begin{array}{ll}1968 . & \text { E. В. Дамаскинский, } \\ & \end{array}$ ПЕРЕСТАНОВОЧНЫХ СООТНОШЕНИЙ ПРЕДСТАВЛЕНИЯ НЕЭКВИВАЛЕНТНЫЕ - представления, наличие которых является характерной особенностью квантовых систем с бесконечным числом степеней свободы. Два представления перестановочных соотношений операторами $R_{\lambda}^{(i)}, i=1,2$, в гильбертовых пространствах $H^{(i)}$ называются неэквивалентными представлениями пере становочных соотношений, если не существует изометрич. отображения $W: H^{(1)} \neg H^{(2)}$ такого, что



\[

R_{\lambda}^{(2)}=W R_{\lambda}^{(1)} W^{-1} \text { для любых } \lambda .

\]



Существование П. с. п. н. было явно отмечено в [1], где они были названы мури отически м и («тьма-тьмущая»). П. с. п. н. имеют разное физич. содержание, так что выбор подходящего представления является первоочередной задачей для строгого описания физич. системы.



Теорема Хаага (см. [2], [3]) означает, что представление Фока соответствует свободным системам и при математически корректном описании взаимодействующих систем требуется привлекать представления неэквивалентные представлению Фока. Одним из следствий наличия П.с. п. н. является сушествование унитарно нереализуемых канонич. преобразований (см. [4], [5]). В соответствии с теоремой Колемана П.с. П. н. необходимы при описании систем со спонтанным нарушением симметрии (см. [6], [7]). В квантовой теории поля П.с. п. н. существенны в решении проблемы инфракрасных расходимостей (см. [8]). В квантовой теории систем многих частиц П.с. П. н. возникают при описании фазовых переходов, сверхтекучести, сверхпроводимости (см. [9]), бозонного и фермионного газов (см. [10], [11]).



См. также Перестановочные соотношения, Перестановочных соотношений представления, Гейзенберга представление.



Jum.: [1] Friedrichs K. O., Mathematical aspects of the quantum theory of fields, N.Y.- L., 1953; [2] Боголюбов H. Н., Логуно в А. А., О к с а к А. И., Т о до ро в И. Т., Общие принципы квантовой теории поля, М., 1987; [3] Э м х Ж., Алгебраические методы в статистической механике и квантовой теории поля, пер. с англ., М., 1976; [4] Бе рез ин Ф.А., Метод вторичного квантования, 2 изд., М., 1986; [5] Завьяло в О. И., Сушко В. Н., в сб.: Статистическая физика и квантовая теория поля, М., 1973, с. 411-40; [6] ГрибА.А., Дамаскинский Е. В., Максимов В. М., «Успехи физич. наук", 1970, т. 102, № 4, с. 587-620; [7] Гри б А. А., Проблемы неинвариантности вакуума в квантовой теории поля, М., 1978; [8] Ку л и ш П. П., Ф а дде е в Л. Д., «Теоретич. и математич. физика", 1970, т. 4, № 2, с. 153-70; [9] H a ag R., «Nuovo Cimento», 1962, v. 25, № 2, p. 287-99; [10] Araki H., Wood s E. J., «J. Math. Phys.», 1963, v. 4, № 5, p. 637-62; [11] Araki H., W yss W., «Helv. Phys. Acta», 1964, v. 37, № 2, p. 136-59. ПЕРЕСТАНОВОЧНЫХ СООТНОШЕНИЙ ПРЕДСТАВЛЕНИЯ УНИТАРНО ЭКВИВАЛЕНТНЫЕ - представления, несущие идентичную физическую информацию и в этом смысле неразличимые (однако для решения конкретных задач целесообразно использование наиболее подходящего из них). Два представления перестановочных соотношений операторами $R_{\lambda}^{(i)}, i=1,2$, в гильбертовых пространствах $H^{(i)}$ называются унитарн о экви вален т ным и, если существует изометрич. отображение $W: H^{(1)} \curvearrowright H^{(2)}$ такое, что



\[

R_{\lambda}^{(2)}=W R_{\lambda}^{(1)} W^{-1} \text { для любых } \lambda .

\]



Согласно теореме единственности Неймана для системы с конечным числом степеней свободы, каждое представление вейлевской формы канонич. коммутационных соотношений унитарно эквивалентно прямой сумме копий Шредингера представления. Наряду с представлением Шредингера в квантовой механике часто используют другие П.с. п. у. э., среди к-рых наиболее известны: представление (чисел заполнения) Фока, метода вторичного квантования, голоморфное представление, гауссово (нормальное) представление, пространство к-рого является $L^{2}$-пространством по гауссовой мере квазиинвариантной относительно трансляций и поэтому допускающее прямое обобщение на случай бесконечных полевых систем. В случае гейзенберговских канонич. коммутационных соотношений даже для конечных систем существуют представления неэквивалентные представлению Шредингера (см. Гейзенберга представление). В случае фермиевских антикоммутационных соотношений для систем с конечным числом степеней свободы согласно теореме Вигнера - Йоордана имеется единственное с точностью до унитарной эквивалентности неприводимое представление. Для системы с бесконечным числом степеней свободы ситуация кардинально меняется и сушествует континуум представлений канонич. коммутационных соотношений и антикоммутационных соотношений, унитарно неэквивалентных стандартному представлению Фока. Эти представления разбиваются на классы унитарно эквивалентных друг другу представлений. Аналогичная ситуация имеется и в случае представлений бесконечномерных групп и алгебр, порождаемых перестановочными соотношениями фундаментальных токов.



Лит. см. при ст. Перестановочных соотношений представления.



Е. В. Дамаскинский.



ПЕРЕХОДА ТОЧКА - см. ФазовЫй переход, Поворота 'точка.



$\mathrm{r} \times \mathrm{i}$





\end{document}